\section{Community artifacts and evaluation agenda}
\label{sec:artifacts}

\subsection{Minimum reproducible automation question}

The first artifact converts recurring clarification requests into a structured technical question. Universal fields cover objective, system and software identity, interface path, method identity, expected behavior, observed behavior, evidence, tests, acceptance criteria, changes since last success, and the requested form of help. Conditional fields cover firmware, licensing, dependencies, layout, resources, geometry, liquid properties, scientific constraints, reproduction scope, and confidentiality. Partial execution activates additional requirements for the physical state after failure and the intervention or recovery sequence.

The full schema is deliberately more precise than a mandatory forum form. A practical interface should present a short universal core and reveal conditional sections based on the question type. Its value can be tested through median clarification turns, time to the first actionable answer, percentage of posts containing version and configuration identity, and actionable-resolution rate. A form that reduces participation or encourages sensitive disclosure would fail its purpose.

\subsection{Resolved-knowledge index}

The second artifact is a governed knowledge record. Each entry contains a stable record identifier, observable problem, bottleneck codes, applicable systems, physical or process context, observed evidence, root-cause status, resolution, validation stage, evidence grade, final disposition where relevant, public artifacts, limitations, supersession lineage, last verification, maintainer, forum provenance, and correction status. Ten seed records demonstrate the structure.

The index treats uncertainty as a state rather than a defect to hide. A record may be active, review due, disputed, superseded, or archived. Root cause may be confirmed, probable, competing, or unknown. Evidence grades prevent a declared value from appearing equivalent to a cross-machine reproduction. This structure addresses the main B1 failure mode by making canonicalization a governed process rather than a static wiki page.

Evaluation should measure repeated-question frequency, new-thread links to existing records, stale-record rate, correction latency, link validity, and the fraction of resolved discussions converted into maintained upstream artifacts. The index would fail if its records accumulate without ownership or if outdated advice receives greater visibility than the original contextual discussion.

\subsection{Device, resource, event, and scheduling registries}

Four further artifacts follow from the technical constructs. A device-integration registry would expose component-level accessibility rather than a single aggregate ranking. A physical-resource registry would bind each field to provenance and evidence grade. A minimal event schema would standardize identifiers and relationships among run, configuration, material, command, acknowledgement, observation, state change, intervention, recovery, and result. A scheduler benchmark suite would publish fully specified scenarios with scientific timing, stability, disposition, and failure constraints.

These artifacts create future denominators. Integration time can be measured per registered interface path. Definition-related faults can be measured per validated resource and run. Recovery outcomes can be measured per incident with known event coverage. Scheduling performance can be measured per fully specified benchmark instance. The community contribution and the research design therefore reinforce each other.
